\section{Conclusion}

This paper introduced the D1 climate exposure dimension of the ROUTE rubric and applied
it to all 227 interstate corridors in the ROUTE v1.2 corpus. Three findings warrant
policy response.

\textbf{Gulf Coast I-10 Louisiana is the most flood-exposed corridor in the national
system.} With 127 consecutive SFHA miles and a mean roadway elevation of 0.8 meters
above sea level, it occupies a qualitatively different risk category than any other
interstate in the corpus. The consecutive-miles metric, not total SFHA miles, is the
correct structural exposure indicator: it captures the evacuation-geometry failure mode
in which an extended flooded span provides no internal escape route.

\textbf{Donner Pass is the highest winter-closure corridor.} Scoring D1 = 7.8 on 900
annual closure-hours (50 events $\times$ 18 hours mean), it demonstrates that the D1
dimension is not synonymous with flood risk. A mountain corridor can be as climate-exposed
as a coastal one through a different mechanism. The D.2 companion paper shows that the
freight cost of Donner closures (\$2.4B/yr) is actually higher than Gulf Coast I-10
(\$1.2B/yr) due to the higher closure frequency and the B1 isolation penalty --- despite
Donner's lower D1 score.

\textbf{The 2050 exposure profile changes the investment priority order.} Under NOAA
Intermediate SLR (0.5m global, 2050), Gulf Coast I-10 Louisiana rises to projected D1 =
9.1 and Gulf Coast Texas to 8.8, while winter corridors increase modestly (+0.2--0.4
points). The correct investment priority is not the 2024 score but the 2050 score, because
infrastructure authorized today will not be operational until the 2035--2040 period at
the earliest.

\textbf{Policy implication.} FHWA's PROTECT program (\$8.7 billion, IIJA \citep{IIJA2021})
currently prioritizes corridors by total SFHA miles. This criterion misses the Gulf Coast
Louisiana corridor's unique structural exposure. Incorporating the consecutive-miles metric
and the 2050 projection into PROTECT eligibility criteria would shift funding priority
toward the corridors where hardening has the highest expected value.

The current PROTECT criteria also do not account for the lead-time urgency: a corridor
that will reach D1 = 9.1 by 2050 requires authorization by 2038 to deliver hardening
infrastructure in time. The program should establish a 2050-exposure tier that triggers
expedited project development timelines for corridors where the projected D1 score exceeds
a threshold (e.g., $\geq 8.0$ under NOAA Intermediate SLR).

The companion paper D.2 \citep{ROUTE_D2} converts these D1 exposure scores into economic
costs of corridor closures, building the full investment case for Gulf Coast I-10 hardening
and the Donner Pass snowshed program. The D1 score is the exposure; D.2 is the consequence.
Together they form the economic basis for Track D investment recommendations in the ROUTE
framework.
